\documentclass[twoside,10pt]{article}

\usepackage[T1]{fontenc}
\usepackage{lmodern}
\usepackage{amsmath,amssymb,amsfonts,amsthm}
\usepackage{mathtools}
\usepackage[small,labelfont=bf,labelsep=period]{caption}
\usepackage[
  letterpaper,
  total={5.5in,8in},
  vmarginratio=1:1,
  hmarginratio=1:1,
  headsep=21pt
]{geometry}
\usepackage{xcolor}
\usepackage[
  breaklinks=false,
  linktocpage=true
]{hyperref}
\hypersetup{
  colorlinks=true,
  linkcolor={purple},
  urlcolor={purple},
  citecolor={blue!80!black},
  pdfauthor={Leslie P. Polzer},
  pdftitle={Algebraic Tunneling for the Restricted Fractional Laplacian on Distant Domains},
  pdfsubject={Leading eigenvalue splitting for distant components of the restricted fractional Laplacian},
  pdfkeywords={fractional Laplacian, eigenvalue splitting, nonlocal double wells, spectral perturbation, shape optimization}
}

\medmuskip=4mu
\thickmuskip=8mu
\setlength{\emergencystretch}{2em}

\numberwithin{equation}{section}

\newtheoremstyle{slbody}
  {8pt}
  {8pt}
  {\normalfont\slshape}
  {}
  {\bfseries}
  {.}
  {0.5em}
  {\thmname{#1}\thmnumber{ #2}%
   \thmnote{ {\normalfont\itshape(\ignorespaces#3\unskip)}}}
\theoremstyle{slbody}
\newtheorem{theorem}{Theorem}[section]
\newtheorem{lemma}[theorem]{Lemma}
\newtheorem{proposition}[theorem]{Proposition}
\newtheorem{corollary}[theorem]{Corollary}
\theoremstyle{definition}
\newtheorem{definition}[theorem]{Definition}
\theoremstyle{remark}
\newtheorem{remark}[theorem]{Remark}

\renewcommand{\qedsymbol}{$\blacksquare$}

\newcommand{\R}{\mathbb{R}}
\newcommand{\Hs}{\widetilde H^s}
\newcommand{\ip}[2]{\langle #1,#2\rangle}
\newcommand{\norm}[1]{\lVert #1\rVert}
\newcommand{\abs}[1]{\lvert #1\rvert}

\title{Algebraic tunneling for the restricted fractional\\
Laplacian on distant domains}
\author{Leslie P. Polzer}
\newcommand{\authoraffiliation}{Independent Researcher}
\newcommand{\authoremail}{polzer@fastmail.com}
\date{September 3, 2026}

\makeatletter
\renewcommand{\section}{%
  \@startsection{section}{1}{\z@}%
    {-3.5ex \@plus -1ex \@minus -.2ex}%
    {2.3ex \@plus .2ex}%
    {\normalfont\large\bfseries}}
\renewcommand{\@seccntformat}[1]{%
  \csname the#1\endcsname.\enspace}

\def\tagform@#1{%
  \maketag@@@{(\ignorespaces{\oldstylenums{#1}}\unskip\@@italiccorr)}}
\renewcommand{\eqref}[1]{%
  \textup{{\normalfont(\oldstylenums{\ref{#1}}}\normalfont)}}

\newcommand{\shorttitle}{ALGEBRAIC TUNNELING ON DISTANT DOMAINS}
\newcommand{\shortauthors}{LESLIE P. POLZER}
\newcommand{\ps@bookheader}{%
  \renewcommand\@oddfoot{\hfil}%
  \renewcommand\@evenfoot{\hfil}%
  \renewcommand\@oddhead{%
    \ifnum\value{page}>1
      {\small\hfil\shortauthors}\hfil\thepage
    \else
      \hfil
    \fi}%
  \renewcommand\@evenhead{%
    \thepage\hfil\small\shorttitle\hfil}}
\pagestyle{bookheader}
\raggedbottom

\renewcommand{\maketitle}{%
  \begin{center}
    {\large\bfseries\@title}\\
    \vskip36pt
    \ifx\@author\@empty
      \mbox{}\par
    \else
      {\scshape\@author}\par
      \medskip
      {\small\authoraffiliation}\par
      \smallskip
      {\small\texttt{\authoremail}}\par
      \smallskip
      {\small\@date}\par
    \fi
    \vskip36pt
  \end{center}}
\makeatother

\renewenvironment{abstract}
  {\quotation\noindent\small{\bfseries\abstractname.\enspace}}
  {\endquotation}

\begin{document}
\maketitle

\begin{abstract}
Let $D\subset B_R(0)\subset\R^d$ be a bounded centrally symmetric Lipschitz
domain, and
let $\Omega_L$ be the union of two translates of $D$, arranged symmetrically
and interchanged by central inversion, whose centers are separated by $L$.  For the
restricted (zero-exterior integral) fractional Laplacian $(-\Delta)^s$,
$0<s<1$, we determine the leading splitting of the two
eigenvalues converging to the one-well ground-state energy.  If $\mu_1$ is
the one-well ground-state eigenvalue, $\phi_1$ is its positive
$L^2$-normalized eigenfunction, and $m_1=\int_D\phi_1$, then
\[
 \lambda_1(\Omega_L)
 =\mu_1-c_{d,s}m_1^2L^{-d-2s}
   +O\!\left(L^{-d-2s-2}+L^{-2d-4s}\right),
\]
\[
 \lambda_2(\Omega_L)
 =\mu_1+c_{d,s}m_1^2L^{-d-2s}
   +O\!\left(L^{-d-2s-2}+L^{-2d-4s}\right).
\]
The proof uses an exact central-inversion reduction to the two operators
$A_D\pm B_L$, an elementary isolated-eigenvalue estimate, and a second-order
Taylor expansion of the interaction kernel.  For two equal balls, the formula
for $\lambda_2$ gives the leading rate in the nonlocal Hong--Krahn--Szeg\H{o}
minimizing sequence.  More generally, central symmetry is unnecessary for
the finite-well result, including its two-well specialization: for any
bounded Lipschitz reference domain and any fixed collection of distinct
translation sites $a_1,\dots,a_N$, the first $N$ eigenvalues are
\[
 \lambda_k=\mu_1+L^{-d-2s}\theta_k
 +O\!\left(L^{-d-2s-2}+L^{-2d-4s}\right),
 \qquad 1\leq k\leq N,
\]
where $\theta_1\leq\cdots\leq\theta_N$ are the eigenvalues of the explicit
matrix with zero diagonal and off-diagonal entries
$-c_{d,s}m_1^2\abs{a_i-a_j}^{-d-2s}$.  Exact cell-integral Galerkin
computations on two and three intervals illustrate the corresponding
fixed-mesh limits, while a mesh study records stabilization of the one-well
interaction coefficient.  In the displayed spectral asymptotics,
$L\to\infty$.
The implicit constants in the two-well formulas depend only on
$d,s,D,R$, and $g$; in the multi-well formula they may also depend on the
fixed $N$ and configuration $(a_1,\dots,a_N)$.

\vskip5pt
\noindent\textbf{Keywords.}\enspace restricted fractional Laplacian;
eigenvalue splitting; nonlocal double wells; spectral perturbation; shape
optimization.

\vskip5pt
\noindent\textbf{MSC2020 Classification.}\enspace 35P15; 35R11.
\end{abstract}

\vskip50pt
\baselineskip=13.5pt
\section{Introduction}

The restricted fractional Laplacian couples disjoint components through its
integral kernel, unlike the local Dirichlet Laplacian, whose spectrum on a
disconnected union is exactly the multiset union of the component spectra.
Thus two distant congruent components exhibit a nonzero splitting that decays
algebraically with their separation.

Brasco and Parini proved that the fractional Hong--Krahn--Szeg\H{o} bound for
the second eigenvalue is strict and that two equal balls form a minimizing
sequence when their mutual distance tends to infinity
\cite[Theorem~6.2]{BrascoParini2016}.  Mixed-operator analogues were proved by
Biagi, Dipierro, Valdinoci, and Vecchi
\cite[Theorem~1.1]{BiagiDipierroValdinociVecchi2023}, and by Goel and Sreenadh
for a local $p$-Laplacian coupled to a compact-kernel nonlocal $p$-Laplacian
\cite[Theorem~1.4]{GoelSreenadh2019}; the latter interaction is finite-range
rather than algebraic.  Parini and Salort placed
such distant components in a general concentration--compactness dichotomy
\cite[Theorems~1.1 and~1.2]{PariniSalort2020}.  For two one-dimensional
patches, L\'eculier and Roquejoffre proved monotonicity and continuity of the
principal eigenvalue throughout the separation range, including contact
\cite[Theorem~1]{LeculierRoquejoffre2023}, but did not identify its leading
large-separation coefficient or the accompanying second-eigenvalue split.
Together these works give qualitative limiting, continuity, and dichotomy
results, but not the leading interaction coefficient for the standard
restricted fractional Laplacian.  We extract that coefficient and the
associated effective finite-dimensional interaction.

The influence of mutual position on the restricted fractional spectrum is
also emphasized by Abatangelo, Felli, and Noris
\cite[Introduction, p.~3]{AbatangeloFelliNoris2020}, in a study of small
fractional-capacity perturbations.  In the shape-optimization setting, Zahl
isolated the exact cross-energy under translations and proposed a signed
point-mass interaction with kernel $\abs{x-y}^{-d-2s}$
\cite[Section~10, equations~(10.3) and~(10.7), and
Conjecture~10.9]{Zahl2026}.  Wu subsequently proved that this discrete
signed-mass model has no local minima, settling that conjecture
\cite[Introduction]{Wu2026}.  In another recent direction, Dipierro, Proietti Lippi,
Sportelli, and Valdinoci studied disconnected domains for generalized
eigenproblems having a superposition of lower-order operators on the
right-hand side
\cite[Theorems~1.4--1.7]{DipierroProiettiLippiSportelliValdinoci2026}.
None of these works derives a large-separation spectral expansion for the
standard restricted fractional Laplacian.  Zahl's discrete energy anticipates
the matrix in \eqref{eq:effective-matrix}; here it is derived from the full
operator, with operator-norm control of the off-diagonal block and spectral-gap
control of the complementary modes, and every ordered cluster eigenvalue is
identified.  We are
not aware of an earlier result giving either the two-component coefficient
with a quantified remainder or this finite-well spectral reduction.

The first result, Theorem~\ref{thm:two-well}, treats the double ground-state
cluster.  Theorem~\ref{thm:multi-well} identifies the entire ground-state
cluster for finitely many translated components with the spectrum of an
explicit weighted interaction matrix; unlike the central-inversion reduction, this
result does not require central symmetry of the reference component.  Its
$N=2$ specialization, Corollary~\ref{cor:symmetry-free-two}, gives the same
leading doublet coefficient without symmetry; symmetry is used only for the
exact scalar-branch decomposition and the matching explicit remainder in
Theorem~\ref{thm:two-well}.  Both proofs are self-contained after the
standard one-well spectral facts.  The~interval provides an explicit
benchmark: one-well eigenvalue and
eigenfunction estimates are available in
\cite[Theorem~1 and Propositions~1--2]{Kwasnicki2012}; sharper eigenvalue
asymptotics and an eigenfunction bound uniform in both the index and the
fractional order are proved in \cite[Theorems~1 and~2]{Zhang2026}.

\section{The one-well operator and exact two-well reduction}

Fix $d\geq1$ and $0<s<1$, and put
\[
 \kappa=d+2s.
\]
For an open set $G\subset\R^d$, let
\[
 \Hs(G)=\{u\in H^s(\R^d):u=0\ \text{a.e. on }\R^d\setminus G\}.
\]
For bounded Lipschitz $G$, this space agrees with the completion of
$C_c^\infty(G)$ in the global $H^s(\R^d)$ norm
\cite[Section~2, equation~(2.1)]{DjitteFallWeth2021}.
The zero-exterior quadratic form is
\begin{equation}\label{eq:form}
 \mathcal q_G[u]
 =\frac{c_{d,s}}2\iint_{\R^d\times\R^d}
   \frac{\abs{u(x)-u(y)}^2}{\abs{x-y}^{\kappa}}\,dx\,dy,
 \qquad u\in\Hs(G),
\end{equation}
where
$c_{d,s}=4^s s\,\Gamma(d/2+s)/(\pi^{d/2}\Gamma(1-s))$
is the standard normalization in the singular-integral definition of
$(-\Delta)^s$.  The form norm
$\bigl(\mathcal q_G[u]+\norm{u}_{L^2(G)}^2\bigr)^{1/2}$
is equivalent to the norm inherited from $H^s(\R^d)$ on $\Hs(G)$.  Thus
\eqref{eq:form} is a densely defined closed form, and we denote its associated
self-adjoint operator on $L^2(G)$ by $A_G$.
We write
$\lambda_1(G)\leq\lambda_2(G)\leq\cdots$ for the eigenvalues of $A_G$,
repeated according to multiplicity.  For an abstract self-adjoint operator
$T$ with compact resolvent, $\lambda_k(T)$ has the analogous meaning.

Let $D\subset B_R(0)$ be a bounded Lipschitz domain satisfying $D=-D$.
The Lipschitz extension property and compactness of the form embedding follow
from \cite[Theorems~5.4 and~7.1]{DiNezzaPalatucciValdinoci2012}.  Hence $A_D$
has compact resolvent.  Write its eigenvalues with multiplicity and denote the
first two distinct levels by
\[
 \mu_1<\mu_2,
 \qquad g=\mu_2-\mu_1>0.
\]
The first eigenvalue is simple and has a strictly positive normalized
eigenfunction by \cite[Theorem~2.8]{BrascoParini2016}, specialized to the
linear case.  Fix that eigenfunction and set
\begin{equation}\label{eq:mass}
 A_D\phi_1=\mu_1\phi_1,
 \qquad \norm{\phi_1}_{L^2(D)}=1,
 \qquad m_1=\int_D\phi_1(x)\,dx>0.
\end{equation}
Simplicity and central-inversion invariance imply $\phi_1(-x)=\phi_1(x)$.

Fix a unit vector $e\in\R^d$.  For $L>2R$, define
\[
 D_L^- =D-\frac L2e,
 \qquad D_L^+=D+\frac L2e,
 \qquad \Omega_L=D_L^-\cup D_L^+.
\]
On $L^2(D)$ define the bounded integral operator
\begin{equation}\label{eq:B}
 (B_Lv)(x)
 =-c_{d,s}\int_D\frac{v(y)}{\abs{Le-x-y}^{\kappa}}\,dy.
\end{equation}
Its kernel is real and symmetric, so $B_L$ is self-adjoint.

\begin{lemma}[Exact central-inversion reduction]\label{lem:block}
The operator $A_{\Omega_L}$ is unitarily equivalent to
\[
 \begin{pmatrix}A_D&B_L\\B_L&A_D\end{pmatrix}
 \quad\text{on }L^2(D)\oplus L^2(D),
\]
and hence to
\begin{equation}\label{eq:direct-sum}
 (A_D+B_L)\oplus(A_D-B_L).
\end{equation}
\end{lemma}

\begin{proof}
For $f\in\Hs(\Omega_L)$ set
\[
 u(x)=f\left(x-\frac L2e\right),
 \qquad
 v(x)=f\left(-x+\frac L2e\right),
 \qquad x\in D.
\]
This defines a unitary map from $L^2(\Omega_L)$ to
$L^2(D)\oplus L^2(D)$.  Translation and central-inversion invariance identify the
two diagonal parts of the form with $\mathcal q_D[u]$ and
$\mathcal q_D[v]$.

To identify the form domains, fix $\chi\in C_c^\infty(\R^d)$ and
$w\in H^s(\R^d)$.  Then
\begin{align*}
 [\chi w]_{H^s(\R^d)}^2
 &\leq 2\norm{\chi}_{L^\infty}^2[w]_{H^s(\R^d)}^2\\
 &\quad+2\int_{\R^d}\abs{w(y)}^2
 \int_{\R^d}\frac{\abs{\chi(x)-\chi(y)}^2}
 {\abs{x-y}^{d+2s}}\,dx\,dy.
\end{align*}
The inner integral is bounded uniformly in $y$: use
$\abs{\chi(x)-\chi(y)}\leq\norm{\nabla\chi}_{L^\infty}\abs{x-y}$
when $\abs{x-y}\leq1$ and
$\abs{\chi(x)-\chi(y)}\leq2\norm{\chi}_{L^\infty}$ otherwise.  The two
radial bounds are integrable because $s<1$ and $s>0$, respectively.
Consequently multiplication by $\chi$ is bounded on $H^s(\R^d)$.

Positive separation permits cutoffs $\chi_\pm$ that equal one on
$D_L^\pm$ and vanish near the other component.  For
$f\in\Hs(\Omega_L)$, the functions $\chi_-f$ and $\chi_+f$ therefore
belong to $H^s(\R^d)$ and, after the above changes of variables, give
$u,v\in\Hs(D)$.  Conversely, translated and reflected zero extensions of
any $u,v\in\Hs(D)$ belong to $H^s(\R^d)$, and their sum is supported in
$\Omega_L$.  Hence the unitary identifies $\Hs(\Omega_L)$ with
$\Hs(D)\oplus\Hs(D)$.

The cross-kernel is bounded by positive separation.  Each separate
zero-exterior form already contains the square interaction with the location
occupied by the other component.  Thus, after the above changes of
variables, expanding the square in \eqref{eq:form} gives exactly
\begin{align*}
 &\mathcal q_{\Omega_L}[f]-\mathcal q_D[u]-\mathcal q_D[v]\\
 &\quad=c_{d,s}\iint_{D\times D}
 \frac{\abs{u(x)-v(y)}^2-\abs{u(x)}^2-\abs{v(y)}^2}
 {\abs{Le-x-y}^{\kappa}}\,dx\,dy\\
 &\quad=-2c_{d,s}\operatorname{Re}\iint_{D\times D}
 \frac{u(x)\overline{v(y)}}{\abs{Le-x-y}^{\kappa}}\,dx\,dy
 =2\operatorname{Re}\ip{u}{B_Lv}.
\end{align*}
The square terms therefore cancel, leaving no additional diagonal
contribution.  This proves the block representation in the sense of closed forms.  Since
$B_L$ is bounded, it also gives the stated operator representation.  The
unitary change of variables
$(u,v)\mapsto 2^{-1/2}(u+v,u-v)$ yields \eqref{eq:direct-sum}.
\end{proof}

\section{The doublet splitting}

We first record the elementary perturbation estimate used for the two scalar
branches.

\begin{lemma}[Ground-state perturbation bound]\label{lem:perturb}
Let $A$ be self-adjoint with compact resolvent, simple lowest eigenvalue
$\mu$, normalized ground state $\phi$, and spectral gap $g>0$.  Let $W$ be
bounded and self-adjoint, put $b=\norm{W}$ and
$w=\ip{\phi}{W\phi}$, and suppose $b\leq g/4$.  Then
\begin{equation}\label{eq:perturb}
 -\frac{2b^2}{g}
 \leq \lambda_1(A+W)-(\mu+w)\leq0.
\end{equation}
\end{lemma}

\begin{proof}
Let $\mathfrak a$ be the closed quadratic form associated with $A$.  The upper
bound follows by testing the Rayleigh quotient with $\phi$.  For a normalized
vector $\psi=a\phi+\eta$ in the form domain of $\mathfrak a$,
where $\eta\perp\phi$ and $t=\norm{\eta}$, the spectral gap gives
\[
 \mathfrak a[\psi]\geq\mu+gt^2.
\]
Since $\abs{a}^2=1-t^2$, self-adjointness gives
\[
 \ip{\psi}{W\psi}-w
 =-t^2w+2\operatorname{Re}\ip{a\phi}{W\eta}+\ip{\eta}{W\eta}.
\]
Using $\abs{w}\leq b$, $\abs{a}\leq1$, and $\norm W=b$,
\begin{align*}
 \mathfrak a[\psi]+\ip{\psi}{W\psi}-(\mu+w)
 &\geq gt^2-2bt^2-2bt\\
 &=(g-2b)t^2-2bt\\
 &= (g-2b)\left(t-\frac{b}{g-2b}\right)^2
   -\frac{b^2}{g-2b}\\
 &\geq-\frac{b^2}{g-2b}
 \geq-\frac{2b^2}{g}.
\end{align*}
Taking the infimum over normalized $\psi$ proves the lower bound.
\end{proof}

The next lemma extracts the interaction coefficient and records why central
symmetry improves the kernel error by one power of $L$.

\begin{lemma}[Interaction-kernel expansion]\label{lem:kernel}
Let
\[
 I_L=\iint_{D\times D}
 \frac{\phi_1(x)\phi_1(y)}{\abs{Le-x-y}^{\kappa}}\,dx\,dy.
\]
For $L\geq4R$,
\begin{equation}\label{eq:kernel-expansion}
 \abs{I_L-m_1^2L^{-\kappa}}
 \leq C_{\kappa,R}m_1^2L^{-\kappa-2},
 \qquad
 C_{\kappa,R}=2^{\kappa+3}\kappa(\kappa+3)R^2.
\end{equation}
Moreover,
\begin{equation}\label{eq:Bnorm}
 \norm{B_L}\leq \beta_L,
 \qquad
 \beta_L=c_{d,s}2^\kappa\abs D\,L^{-\kappa}.
\end{equation}
\end{lemma}

\begin{proof}
Set $F_L(z)=\abs{Le-z}^{-\kappa}$.  For $\abs z\leq2R$ and $L\geq4R$,
the segment from $0$ to $z$ remains at distance at least $L/2$ from $Le$.
Direct differentiation gives
\[
 \norm{D^2F_L(z)}
 \leq \kappa(\kappa+3)\abs{Le-z}^{-\kappa-2}.
\]
Taylor's theorem therefore yields
\[
 \abs{F_L(z)-L^{-\kappa}
 -\kappa L^{-\kappa-1}e\cdot z}
 \leq C_{\kappa,R}L^{-\kappa-2}
\]
for $\abs z\leq2R$, with the displayed value of $C_{\kappa,R}$.
Because $\phi_1$ is even,
\[
 \int_Dx\phi_1(x)\,dx=0.
\]
Consequently the integral of the linear Taylor term with $z=x+y$ vanishes.
Since $\phi_1$ is positive, integration of the remainder gives
\eqref{eq:kernel-expansion}.

Finally, $\abs{Le-x-y}\geq L-2R\geq L/2$.  The Schur test applied to
\eqref{eq:B} gives
\[
 \norm{B_L}
 \leq c_{d,s}\abs D(L-2R)^{-\kappa}
 \leq c_{d,s}2^\kappa\abs D\,L^{-\kappa},
\]
which is \eqref{eq:Bnorm}.
\end{proof}

\begin{theorem}[Two-well algebraic splitting]\label{thm:two-well}
Let $D\subset B_R(0)$ be a bounded centrally symmetric Lipschitz domain, and
let $\Omega_L$ be defined above.  There is $L_0<\infty$, determined by
$d,s,D,R$ and the one-well gap $g$, such that for $L\geq L_0$ the first two
eigenvalues of $A_{\Omega_L}$ satisfy
\begin{align}
 \abs{\lambda_1(\Omega_L)
  -(\mu_1-c_{d,s}m_1^2L^{-\kappa})}
 &\leq c_{d,s}C_{\kappa,R}m_1^2L^{-\kappa-2}
       +\frac{2\beta_L^2}{g},\label{eq:lambda1}\\
 \abs{\lambda_2(\Omega_L)
  -(\mu_1+c_{d,s}m_1^2L^{-\kappa})}
 &\leq c_{d,s}C_{\kappa,R}m_1^2L^{-\kappa-2}
       +\frac{2\beta_L^2}{g}.\label{eq:lambda2}
\end{align}
In particular, as $L\to\infty$,
\begin{equation}\label{eq:gap-limit}
 L^{d+2s}\bigl(\lambda_2(\Omega_L)-\lambda_1(\Omega_L)\bigr)
 \longrightarrow 2c_{d,s}m_1^2.
\end{equation}
\end{theorem}

\begin{proof}
By Lemma~\ref{lem:block}, the spectrum is the union of the spectra of
$A_D+B_L$ and $A_D-B_L$.  Put
\[
 t_L=\ip{\phi_1}{B_L\phi_1}=-c_{d,s}I_L<0.
\]
Choose $L_0\geq4R$ so large that, for $L\geq L_0$,
\begin{equation}\label{eq:L0-conditions}
 \beta_L\leq\frac g4,
 \qquad
 2c_{d,s}m_1^2(L+2R)^{-\kappa}>\frac{4\beta_L^2}{g}.
\end{equation}
Such a choice exists because the left-hand side of the second inequality is
of order $L^{-\kappa}$ and its right-hand side is of order $L^{-2\kappa}$.

Apply Lemma~\ref{lem:perturb} first with $W=B_L$ and then with $W=-B_L$.
The two branch ground-state energies have the form
\begin{equation}\label{eq:branches}
 E_+(L)=\mu_1+t_L+\rho_+(L),
 \qquad
 E_-(L)=\mu_1-t_L+\rho_-(L),
\end{equation}
where
\[
 \abs{\rho_\pm(L)}\leq\frac{2\beta_L^2}{g}.
\]
The second eigenvalue of either branch is at least
$\mu_2-\beta_L\geq\mu_1+3g/4$, whereas both branch ground-state energies are
at most $\mu_1+\beta_L\leq\mu_1+g/4$.  Thus the first two eigenvalues of the
full operator are the two numbers in \eqref{eq:branches}.

Since $\abs{Le-x-y}\leq L+2R$ and $\phi_1>0$,
\[
 I_L\geq m_1^2(L+2R)^{-\kappa}.
\]
The second condition in \eqref{eq:L0-conditions} and
\eqref{eq:branches} imply $E_+(L)<E_-(L)$.  Hence
$\lambda_1(\Omega_L)=E_+(L)$ and
$\lambda_2(\Omega_L)=E_-(L)$.  Equations \eqref{eq:lambda1} and
\eqref{eq:lambda2} now follow from Lemma~\ref{lem:kernel}.  Subtracting the
two expansions and multiplying by $L^\kappa$ proves
\eqref{eq:gap-limit}, since $\kappa>0$.
\end{proof}

\begin{corollary}[Rate for the distant-ball minimizing sequence]
Let $D=B_R(0)$ and let $\Omega_L$ be the union of two radius-$R$ balls whose
centers are separated by $L$, so that $\abs{\Omega_L}=2\abs{B_R}$.  Then, as
$L\to\infty$,
\[
 L^{d+2s}\bigl(\lambda_2(\Omega_L)-\lambda_1(B_R)\bigr)
 \longrightarrow c_{d,s}
 \left(\int_{B_R}\phi_1(x)\,dx\right)^2.
\]
Thus the minimizing sequence in
\cite[Theorem~6.2]{BrascoParini2016} approaches its limiting value at the
explicit leading order $L^{-d-2s}$.
\end{corollary}

\begin{proof}
This is \eqref{eq:lambda2} with $D=B_R(0)$ and
$\mu_1=\lambda_1(B_R)$.
\end{proof}

\section{Finitely many wells}\label{sec:multi-well}

The two-well central-inversion decomposition is special to a pair.  For finitely
many components we instead keep the full block operator and compare its
lowest spectral cluster with its compression to the translated one-well
ground states.

Throughout this section, let $D\subset B_R(0)$ be any bounded Lipschitz
domain; central symmetry is not assumed.  The compactness and ground-state
results \cite[Theorems~5.4 and~7.1]{DiNezzaPalatucciValdinoci2012} and
\cite[Theorem~2.8]{BrascoParini2016} apply unchanged.  Write
$\mu_1<\mu_2$, $g=\mu_2-\mu_1>0$, choose the positive normalized ground
state $A_D\phi_1=\mu_1\phi_1$, and set
$m_1=\int_D\phi_1(x)\,dx>0$.

\begin{definition}[Multi-well geometry]\label{def:multi-geometry}
Fix an integer $N\geq2$ and distinct points
$\mathbf a=(a_1,\dots,a_N)$ in $\R^d$.  Set
\begin{equation}\label{eq:geometry-constants}
 \delta_{\mathbf a}=\min_{i\ne j}\abs{a_i-a_j},
 \qquad
 \Sigma_p(\mathbf a)
 =\max_{1\leq i\leq N}\sum_{j\ne i}\abs{a_i-a_j}^{-p}
 \quad (p>0).
\end{equation}
For $L\delta_{\mathbf a}>2R$, define
\begin{equation}\label{eq:multi-domain}
 D_{j,L}=D+La_j,
 \qquad
 \Omega_{\mathbf a,L}=\bigcup_{j=1}^N D_{j,L}.
\end{equation}
The sets in this union are pairwise disjoint.
\end{definition}

On $\mathcal H_N=\bigoplus_{j=1}^N L^2(D)$, let
$\mathcal A_0=\bigoplus_{j=1}^N A_D$.  For $i\ne j$ define
\begin{equation}\label{eq:multi-block}
 (B_{ij,L}v)(x)
 =-c_{d,s}\int_D
 \frac{v(y)}{\abs{L(a_i-a_j)+x-y}^{\kappa}}\,dy,
 \qquad B_{ii,L}=0,
\end{equation}
and let $\mathcal V_L=(B_{ij,L})_{i,j=1}^N$.

\begin{lemma}[Exact multi-well reduction and norm bound]
\label{lem:multi-block}
If $L\delta_{\mathbf a}\geq4R$, then $A_{\Omega_{\mathbf a,L}}$ is
unitarily equivalent to $\mathcal A_0+\mathcal V_L$.  Moreover,
\begin{equation}\label{eq:gamma-L}
 \norm{\mathcal V_L}\leq\gamma_L,
 \qquad
 \gamma_L
 =c_{d,s}2^\kappa\abs D\,\Sigma_\kappa(\mathbf a)L^{-\kappa}.
\end{equation}
\end{lemma}

\begin{proof}
For $f\in L^2(\Omega_{\mathbf a,L})$, put
$u_j(x)=f(x+La_j)$ for $x\in D$.  This is a unitary identification with
$\mathcal H_N$.  Translation invariance gives the diagonal operator
$\mathcal A_0$.  The cutoff argument in the proof of Lemma~\ref{lem:block}
identifies the form domain with $\bigoplus_{j=1}^N\Hs(D)$, and positive
separation makes every cross-kernel bounded.  For $i<j$, direct expansion of
the form
\eqref{eq:form} gives the cross term
\[
 -2c_{d,s}\operatorname{Re}\iint_{D\times D}
 \frac{u_i(x)\overline{u_j(y)}}
 {\abs{L(a_i-a_j)+x-y}^{\kappa}}\,dx\,dy.
\]
Since $B_{ji,L}=B_{ij,L}^*$, summing these terms proves the exact block
representation in the sense of closed forms and hence as operators.

For $i\ne j$ and $x,y\in D$,
\[
 \abs{L(a_i-a_j)+x-y}
 \geq L\abs{a_i-a_j}-2R
 \geq\frac L2\abs{a_i-a_j}.
\]
The Schur test therefore gives
\begin{equation}\label{eq:multi-block-norm}
 \norm{B_{ij,L}}
 \leq c_{d,s}2^\kappa\abs D\,
 L^{-\kappa}\abs{a_i-a_j}^{-\kappa}.
\end{equation}
For $u=(u_1,\dots,u_N)$, set $r_j=\norm{u_j}_{L^2(D)}$ and let
$G$ be the symmetric matrix with entries
$G_{ij}=\norm{B_{ij,L}}$.  Then
\[
 \norm{\mathcal V_Lu}_{\mathcal H_N}
 \leq\norm{Gr}_{\ell^2}
 \leq\left(\max_i\sum_jG_{ij}\right)\norm r_{\ell^2}.
\]
Combining this inequality with \eqref{eq:multi-block-norm} proves
\eqref{eq:gamma-L}.
\end{proof}

We next record the finite-multiplicity perturbation estimate that replaces
Lemma~\ref{lem:perturb}.

\begin{lemma}[Isolated spectral-cluster bound]\label{lem:cluster}
Let $A$ be self-adjoint with compact resolvent.  Suppose its lowest
eigenvalue $\mu$ has multiplicity $N$, let $P$ be the corresponding spectral
projection, and suppose
\[
 A\big|_{P^\perp}\geq\mu+g
\]
in the quadratic-form sense for some $g>0$.  Let $W$ be bounded and
self-adjoint, set $b=\norm W$, and assume $b\leq g/4$.  If
$\eta_1\leq\cdots\leq\eta_N$ are the eigenvalues of $PWP$ on
$\operatorname{ran}P$, then, for $1\leq k\leq N$,
\begin{equation}\label{eq:cluster-bound}
 -\frac{2b^2}{g}
 \leq\lambda_k(A+W)-(\mu+\eta_k)\leq0.
\end{equation}
In addition,
\begin{equation}\label{eq:cluster-separation}
 \lambda_{N+1}(A+W)\geq\mu+g-b\geq\mu+\frac{3g}{4}.
\end{equation}
\end{lemma}

\begin{proof}
The min--max principle applied to the span of the first $k$ eigenvectors of
$PWP$ gives the upper bound in \eqref{eq:cluster-bound}.  For the lower
bound, write $\psi=p+q$ with $p=P\psi$ and $q=(I-P)\psi$.  Then
\begin{align*}
 \ip{\psi}{(A+W)\psi}
 &\geq\mu\norm\psi^2+\ip{p}{PWPp}
 +(g-b)\norm q^2-2b\norm p\norm q\\
 &\geq\mu\norm\psi^2+\ip{p}{PWPp}
 -\frac{2b^2}{g}\norm p^2
 +\left(\frac g2-b\right)\norm q^2.
\end{align*}
The second line uses
$2b\norm p\norm q\leq2b^2g^{-1}\norm p^2+(g/2)\norm q^2$.
Thus, in the form sense, $A+W$ dominates the direct sum of
\[
 \mu I+PWP-\frac{2b^2}{g}I
 \quad\text{on }\operatorname{ran}P
\]
and $(\mu+g/2-b)I$ on $P^\perp$.  Since
$\eta_N\leq b\leq g/4$, we have
\[
 \eta_N-\frac{2b^2}{g}
 \leq b-\frac{2b^2}{g}<\frac g2-b.
\]
The strict inequality follows because the difference between the last two
quantities is $(g-2b)^2/(2g)>0$.
Hence the first $N$ eigenvalues of this comparison operator are the $N$
eigenvalues of its $P$-block.  Another
application of the min--max principle gives the lower bound in
\eqref{eq:cluster-bound}.  Finally, $A+W\geq A-bI$ and the min--max principle
give \eqref{eq:cluster-separation}.
\end{proof}

Let $P_1$ be the orthogonal projection in $\mathcal H_N$ onto the span of
the $N$ vectors whose $j$th component is $\phi_1$ and whose other components
vanish.  In this basis the compression $P_1\mathcal V_LP_1$ is represented
by the real symmetric matrix $T_L=(t_{ij}(L))$, where
\begin{equation}\label{eq:exact-compression}
 t_{ii}(L)=0,
 \qquad
 t_{ij}(L)=-c_{d,s}\iint_{D\times D}
 \frac{\phi_1(x)\phi_1(y)}
 {\abs{L(a_i-a_j)+x-y}^{\kappa}}\,dx\,dy
 \quad(i\ne j).
\end{equation}

\begin{definition}[Effective interaction matrix]\label{def:effective-matrix}
Define the real symmetric $N\times N$ matrix $M_{\mathbf a}$ by
\begin{equation}\label{eq:effective-matrix}
 (M_{\mathbf a})_{ii}=0,
 \qquad
 (M_{\mathbf a})_{ij}
 =-c_{d,s}m_1^2\abs{a_i-a_j}^{-\kappa}
 \quad(i\ne j).
\end{equation}
Write its ordered eigenvalues as
$\theta_1\leq\cdots\leq\theta_N$.
\end{definition}

\begin{lemma}[Compression error]\label{lem:compression}
If $L\delta_{\mathbf a}\geq4R$, then
\begin{equation}\label{eq:compression-error}
 \norm{T_L-L^{-\kappa}M_{\mathbf a}}_{\ell^2\to\ell^2}
 \leq\varepsilon_L,
 \qquad
 \varepsilon_L
 =c_{d,s}C_{\kappa,R}m_1^2
 \Sigma_{\kappa+2}(\mathbf a)L^{-\kappa-2}.
\end{equation}
\end{lemma}

\begin{proof}
Fix $i\ne j$, put $a=a_i-a_j$, and set
$F(z)=\abs{La+z}^{-\kappa}$.  For $\abs z\leq2R$, the segment from $0$ to
$z$ stays at distance at least $L\abs a/2$ from $-La$.  As in the proof of
Lemma~\ref{lem:kernel}, Taylor's theorem gives
\begin{equation}\label{eq:multi-taylor}
 \abs{F(z)-F(0)-DF(0)z}
 \leq C_{\kappa,R}L^{-\kappa-2}\abs a^{-\kappa-2}.
\end{equation}
The linear term vanishes after integration because
\[
 \iint_{D\times D}\phi_1(x)\phi_1(y)(x-y)\,dx\,dy=0.
\]
This cancellation is algebraic and does not require central symmetry of
$D$.  Using $F(0)=L^{-\kappa}\abs a^{-\kappa}$ in
\eqref{eq:exact-compression} and integrating \eqref{eq:multi-taylor} gives
\[
 \abs{t_{ij}(L)-L^{-\kappa}(M_{\mathbf a})_{ij}}
 \leq c_{d,s}C_{\kappa,R}m_1^2
 L^{-\kappa-2}\abs{a_i-a_j}^{-\kappa-2}.
\]
The operator norm of a real symmetric matrix is at most its largest absolute
row sum.  Taking that row sum proves \eqref{eq:compression-error}.
\end{proof}

\begin{theorem}[Finite multi-well effective Hamiltonian]
\label{thm:multi-well}
Let $D$, $\mu_1$, $g$, $\phi_1$, and $m_1$ be as fixed at the beginning of
this section, and let the multi-well geometry be as in
Definition~\ref{def:multi-geometry}.  If
\begin{equation}\label{eq:multi-size-conditions}
 L\delta_{\mathbf a}\geq4R,
 \qquad
 \gamma_L\leq\frac g4,
\end{equation}
then, for $1\leq k\leq N$,
\begin{equation}\label{eq:multi-eigenvalue-bound}
 \abs{\lambda_k(\Omega_{\mathbf a,L})
 -(\mu_1+L^{-\kappa}\theta_k)}
 \leq\varepsilon_L+\frac{2\gamma_L^2}{g}.
\end{equation}
The upper edge of the cluster and the next eigenvalue satisfy
\begin{equation}\label{eq:multi-cluster-edges}
 \lambda_N(\Omega_{\mathbf a,L})
 \leq\mu_1+\gamma_L\leq\mu_1+\frac g4,
 \qquad
 \lambda_{N+1}(\Omega_{\mathbf a,L})
 \geq\mu_1+g-\gamma_L\geq\mu_1+\frac{3g}{4}.
\end{equation}
In particular, the cluster gap has the quantitative lower bound
\begin{equation}\label{eq:multi-cluster-gap}
 \lambda_{N+1}(\Omega_{\mathbf a,L})-
 \lambda_N(\Omega_{\mathbf a,L})
 \geq g-2\gamma_L\geq\frac g2.
\end{equation}
Consequently, as $L\to\infty$, for every $1\leq k\leq N$,
\begin{equation}\label{eq:multi-limit}
 L^{d+2s}\bigl(\lambda_k(\Omega_{\mathbf a,L})-\mu_1\bigr)
 \longrightarrow\theta_k.
\end{equation}
\end{theorem}

\begin{proof}
Lemma~\ref{lem:multi-block} identifies the full operator with
$\mathcal A_0+\mathcal V_L$.  The lowest eigenvalue of $\mathcal A_0$ is
$\mu_1$, with multiplicity $N$, its spectral projection is $P_1$, and the
rest of its spectrum lies in $[\mu_1+g,\infty)$.  Apply
Lemma~\ref{lem:cluster} with $A=\mathcal A_0$, $W=\mathcal V_L$, and
$b=\norm{\mathcal V_L}\leq\gamma_L$.  If
$\tau_1(L)\leq\cdots\leq\tau_N(L)$ are the eigenvalues of $T_L$, this gives
\begin{equation}\label{eq:exact-cluster-comparison}
 -\frac{2\gamma_L^2}{g}
 \leq\lambda_k(\Omega_{\mathbf a,L})-(\mu_1+\tau_k(L))\leq0.
\end{equation}
The finite-dimensional min--max principle and
Lemma~\ref{lem:compression} give
\[
 \abs{\tau_k(L)-L^{-\kappa}\theta_k}\leq\varepsilon_L.
\]
Combining these two estimates proves \eqref{eq:multi-eigenvalue-bound}, and
the upper bound in \eqref{eq:exact-cluster-comparison}, together with
$\tau_N(L)\leq\norm{T_L}\leq\norm{\mathcal V_L}\leq\gamma_L$, gives the
upper-edge estimate in \eqref{eq:multi-cluster-edges}.  The lower bound for
$\lambda_{N+1}$ follows from \eqref{eq:cluster-separation}.  Subtracting the
two estimates proves \eqref{eq:multi-cluster-gap}.
Finally, $L^\kappa\varepsilon_L=O(L^{-2})$ and
$L^\kappa\gamma_L^2=O(L^{-\kappa})$; since $\kappa=d+2s>0$, multiplying
\eqref{eq:multi-eigenvalue-bound} by $L^\kappa$ proves
\eqref{eq:multi-limit}.
\end{proof}

\begin{corollary}[Symmetry-free two-well asymptotic]
\label{cor:symmetry-free-two}
Let $D\subset B_R(0)$ be any bounded Lipschitz domain, let $a_1\ne a_2$, and
put $r=\abs{a_1-a_2}$.  For
$\Omega_L=(D+La_1)\cup(D+La_2)$, as $L\to\infty$,
\begin{align*}
 \lambda_1(\Omega_L)
 &=\mu_1-c_{d,s}m_1^2r^{-\kappa}L^{-\kappa}
   +O\!\left(L^{-\kappa-2}+L^{-2\kappa}\right),\\
 \lambda_2(\Omega_L)
 &=\mu_1+c_{d,s}m_1^2r^{-\kappa}L^{-\kappa}
   +O\!\left(L^{-\kappa-2}+L^{-2\kappa}\right).
\end{align*}
In particular,
\[
 L^{d+2s}\bigl(\lambda_2(\Omega_L)-\lambda_1(\Omega_L)\bigr)
 \longrightarrow 2c_{d,s}m_1^2r^{-d-2s}.
\]
\end{corollary}

\begin{proof}
For $N=2$, the effective matrix is
\[
 M_{\mathbf a}
 =\begin{pmatrix}
 0&-c_{d,s}m_1^2r^{-\kappa}\\
 -c_{d,s}m_1^2r^{-\kappa}&0
 \end{pmatrix},
\]
whose ordered eigenvalues are
$-c_{d,s}m_1^2r^{-\kappa}$ and $c_{d,s}m_1^2r^{-\kappa}$.
Theorem~\ref{thm:multi-well} gives both expansions because
$\varepsilon_L=O(L^{-\kappa-2})$ and
$\gamma_L^2=O(L^{-2\kappa})$.  Subtracting them and using $\kappa=d+2s$
gives the limit.
\end{proof}

\begin{corollary}[Recovery of the doublet]\label{cor:multi-two}
Assume in addition that $D=-D$.  For $N=2$, $a_1=-e/2$, and $a_2=e/2$,
the right-hand side of \eqref{eq:multi-eigenvalue-bound}, including its
explicit constants, agrees with the bounds
\eqref{eq:lambda1}--\eqref{eq:lambda2}.
\end{corollary}

\begin{proof}
In this case
\[
 M_{\mathbf a}
 =\begin{pmatrix}0&-c_{d,s}m_1^2\\-c_{d,s}m_1^2&0\end{pmatrix},
\]
whose ordered eigenvalues are the two stated numbers.  Moreover,
$\Sigma_\kappa(\mathbf a)=\Sigma_{\kappa+2}(\mathbf a)=1$, so
$\gamma_L=\beta_L$.  In the exact compression, central symmetry of $D$ and
evenness of $\phi_1$ show, after the change of variables $y\mapsto-y$, that
the translated two-well kernel integral equals $I_L$.  Thus the compression
errors, and hence the two displayed remainder bounds, coincide.
\end{proof}

\begin{corollary}[Collective ground state]\label{cor:collective-ground}
For every multi-well configuration, $\theta_1<0$ is simple and admits an
eigenvector with strictly positive entries, while $\theta_N>0$.
For every $L$ admitted by Definition~\ref{def:multi-geometry},
$\lambda_1(\Omega_{\mathbf a,L})$ is simple.  Moreover, for all sufficiently
large $L$,
$\lambda_1(\Omega_{\mathbf a,L})<\mu_1<\lambda_N(\Omega_{\mathbf a,L})$.
\end{corollary}

\begin{proof}
The off-diagonal entries of $M_{\mathbf a}$ are strictly negative, and its
Rayleigh quotient at $(1,\dots,1)$ is negative, so $\theta_1<0$.  Replacing a
Rayleigh minimizer $z$ by its componentwise absolute value cannot increase
the quotient; equality forces its nonzero components to have the same sign,
and the eigenvalue equation rules out a zero component.  Thus a ground-state
eigenvector has one strict sign.  If the ground eigenspace had dimension at
least two, it would contain a nonzero vector with a zero component, a
contradiction.  Hence $\theta_1$ is simple.  Finally,
$\operatorname{tr}M_{\mathbf a}=0$ gives $\theta_N>0$.  Put
$r_L=\varepsilon_L+2\gamma_L^2/g$.  Theorem~\ref{thm:multi-well} gives
$L^\kappa r_L\to0$.  Simplicity of
$\lambda_1(\Omega_{\mathbf a,L})$ for every admissible $L$ follows from
\cite[Theorem~2.8]{BrascoParini2016}, applied to the bounded open set
$\Omega_{\mathbf a,L}$.  Choose $L$ so large that
\[
 L^\kappa r_L
 <\min\left\{-\theta_1,\theta_N\right\}.
\]
Then \eqref{eq:multi-eigenvalue-bound} gives
$\lambda_1<\mu_1<\lambda_N$, which proves the remaining assertions.
\end{proof}

\begin{corollary}[Regular-simplex cluster]\label{cor:simplex}
Suppose the centers satisfy $\abs{a_i-a_j}=r$ for all $i\ne j$, as for the
vertices of a regular simplex.  Then $N\leq d+1$.  Put
\[
 w=c_{d,s}m_1^2r^{-\kappa}.
\]
Then $\theta_1=-(N-1)w$ and
$\theta_2=\cdots=\theta_N=w$.  Thus one collective level is lowered by
$(N-1)wL^{-\kappa}$, while the remaining $N-1$ levels are raised by
$wL^{-\kappa}$, up to the remainder in
\eqref{eq:multi-eigenvalue-bound}.
\end{corollary}

\begin{proof}
After translating $a_N$ to the origin, the $N-1$ vectors
$a_i-a_N$ have Gram matrix with diagonal entries $r^2$ and off-diagonal
entries $r^2/2$.  This matrix is positive definite, so $N-1\leq d$.
Here $M_{\mathbf a}=-w(J-I)$, where $J$ is the all-ones matrix.  The vector
$(1,\dots,1)$ has eigenvalue $-(N-1)w$, and its orthogonal complement has
eigenvalue $w$.
\end{proof}

\section{A reproducible interval check}\label{sec:numerics}

We finish with a conforming fixed-mesh Galerkin check of the two- and
three-well limits.  It tests the formulas and is not an input to their proofs.
No $h\to0$ discretization-error estimate is claimed; the calculation tests
the proved $L\to\infty$ asymptotic at fixed mesh.  Take $d=1$, $0<s<1/2$,
and partition each unit interval into cells of length $h$.  The cell
indicators have finite form energy by Proposition~\ref{prop:cell-matrix}, so
they span a conforming finite-dimensional subspace of the form domain.

\begin{proposition}[Exact cell-integral matrix]\label{prop:cell-matrix}
Let $I_1,\dots,I_M$ be pairwise disjoint cells of common length $h$, with
centers $x_\alpha$, and set $G=\bigcup_{\alpha=1}^M I_\alpha$,
$r_{\alpha\beta}=\abs{x_\alpha-x_\beta}$, and $q=1-2s$.  In the
$L^2(G)$-orthonormal basis $e_\alpha=h^{-1/2}\mathbf 1_{I_\alpha}$, the
Galerkin matrix $H_h$ of the zero-exterior form $\mathcal q_G$ has entries
\begin{align}
 (H_h)_{\alpha\alpha}
 &=\frac{c_{1,s}}{h}\frac{2h^q}{2s q},
 \label{eq:cell-diagonal}\\
 (H_h)_{\alpha\beta}
 &=-\frac{c_{1,s}}{h}
 \frac{2r_{\alpha\beta}^q-(r_{\alpha\beta}+h)^q
 -(r_{\alpha\beta}-h)^q}{2s q},
 \qquad \alpha\ne\beta.
 \label{eq:cell-offdiagonal}
\end{align}
\end{proposition}

\begin{proof}
For two disjoint cells with center distance $r\geq h$, two direct
integrations give
\[
 \iint_{I_\alpha\times I_\beta}\abs{x-y}^{-1-2s}\,dx\,dy
 =\frac{2r^q-(r+h)^q-(r-h)^q}{2s q}.
\]
The off-diagonal form entry is minus $c_{1,s}$ times this integral.  For one
cell,
\[
 \int_{I_\alpha}\int_{\R\setminus I_\alpha}
 \abs{x-y}^{-1-2s}\,dy\,dx
 =\frac{2h^q}{2s q}.
\]
Dividing both expressions by $h$ for the normalized indicators proves
\eqref{eq:cell-diagonal}--\eqref{eq:cell-offdiagonal}.
\end{proof}

\begin{corollary}[Fixed-mesh cluster asymptotic]
\label{cor:discrete-cluster}
Let $D=(-1/2,1/2)$, and fix a uniform partition of $D$.  Let $\mu_{1,h}$
and $\phi_{1,h}$ be the lowest eigenvalue and a positive,
$L^2$-normalized piecewise-constant eigenfunction of the one-well Galerkin
matrix.  Put $m_{1,h}=\int_D\phi_{1,h}$.  Translate this same partition to
each component $D+La_j$, and write the ordered eigenvalues of the resulting
Galerkin matrix on $\bigcup_{j=1}^N(D+La_j)$ as $\lambda_{k,h}(L)$.  Then,
as $L\to\infty$, for fixed distinct
$a_1,\dots,a_N$,
\begin{equation}\label{eq:discrete-cluster}
 \lambda_{k,h}(L)
 =\mu_{1,h}+L^{-1-2s}\theta_{k,h}
 +O_{h,\mathbf a,s}\!\left(L^{-3-2s}+L^{-2-4s}\right),
 \qquad 1\leq k\leq N,
\end{equation}
where $\theta_{1,h}\leq\cdots\leq\theta_{N,h}$ are the eigenvalues of the
matrix with zero diagonal and off-diagonal entries
$-c_{1,s}m_{1,h}^2\abs{a_i-a_j}^{-1-2s}$.
\end{corollary}

\begin{proof}
All off-diagonal entries of the one-well matrix in
\eqref{eq:cell-offdiagonal} are negative.  The Rayleigh-quotient argument in
the proof of Corollary~\ref{cor:collective-ground} therefore gives a simple
lowest eigenvalue and a strictly positive eigenvector.  On the direct sum of
$N$ fixed Galerkin spaces, the
block expansion used in Lemma~\ref{lem:multi-block} is exact, its
off-diagonal norm is $O_{h,\mathbf a,s}(L^{-1-2s})$, and the complement of
the $N$ copied ground states is separated by the positive discrete
one-well gap.  The Taylor argument in Lemma~\ref{lem:compression}, now in
Euclidean operator norm, gives an
$O_{h,\mathbf a,s}(L^{-3-2s})$ compression error.  Applying
Lemma~\ref{lem:cluster} to this finite matrix gives the additional
$O_{h,\mathbf a,s}(L^{-2-4s})$ term and proves
\eqref{eq:discrete-cluster}.
\end{proof}

For the computation we take $s=1/4$, $D=(-1/2,1/2)$, and the standard
normalization
\[
 c_{1,s}=\frac{4^s s\,\Gamma(1/2+s)}
 {\sqrt\pi\,\Gamma(1-s)}.
\]
Table~\ref{tab:mesh} records the one-well quantities.  The stabilization of
$m_{1,h}$ is the relevant check for the interaction coefficient.

\begin{table}[ht]
\centering
\begin{tabular}{rccc}
\hline
cells & $\mu_{1,h}$ & $\mu_{2,h}-\mu_{1,h}$ & $m_{1,h}$\\ \hline
60  & 1.3794549573 & 0.9105539931 & 0.9628386433\\
120 & 1.3753692130 & 0.8998883318 & 0.9633505837\\
180 & 1.3741292905 & 0.8969647744 & 0.9636021778\\
\hline
\end{tabular}
\caption{One-well Galerkin data for $s=1/4$.}
\label{tab:mesh}
\end{table}

On the 180-cell mesh, the two-well component centers are $(-L/2,L/2)$,
while the three-well centers are $(-L,0,L)$.  The corresponding fixed-mesh
effective spectra are
\begin{align*}
 (-L/2,L/2):&\quad (-0.18521477,0.18521477),\\
 (-L,0,L):&\quad
 (-0.29671332,0.06548331,0.23123001).
\end{align*}
The direct Galerkin eigenvalues are shown in Table~\ref{tab:clusters}; the
last column is the largest absolute difference from the corresponding
effective-matrix eigenvalue.  These errors are computed from the
full-precision values before the shifts and effective spectra are rounded for
display.

\begin{table}[ht]
\centering
\begin{tabular}{rccc}
\hline
$L$ & $L^{3/2}(\lambda_{1,h}-\mu_{1,h})$
    & $L^{3/2}(\lambda_{2,h}-\mu_{1,h})$ & max. error\\ \hline
3  & $-0.190740$ & $ 0.189832$ & $5.52\times10^{-3}$\\
6  & $-0.186544$ & $ 0.186335$ & $1.33\times10^{-3}$\\
12 & $-0.185551$ & $ 0.185486$ & $3.36\times10^{-4}$\\
24 & $-0.185286$ & $ 0.185295$ & $8.04\times10^{-5}$\\
\hline
\end{tabular}

\vskip10pt
\begin{tabular}{rcccc}
\hline
$L$ & $L^{3/2}(\lambda_{1,h}-\mu_{1,h})$
    & $L^{3/2}(\lambda_{2,h}-\mu_{1,h})$
    & $L^{3/2}(\lambda_{3,h}-\mu_{1,h})$ & max. error\\ \hline
3  & $-0.304984$ & $0.065562$ & $0.237544$ & $8.27\times10^{-3}$\\
6  & $-0.298738$ & $0.065554$ & $0.232747$ & $2.03\times10^{-3}$\\
12 & $-0.297232$ & $0.065510$ & $0.231597$ & $5.19\times10^{-4}$\\
24 & $-0.296843$ & $0.065504$ & $0.231323$ & $1.29\times10^{-4}$\\
\hline
\end{tabular}
\caption{Scaled two-well (top) and three-well (bottom) cluster shifts.}
\label{tab:clusters}
\end{table}

The scaled errors decrease to zero in both configurations, as asserted at
fixed mesh by Corollary~\ref{cor:discrete-cluster}.  The calculation uses the
exact entries \eqref{eq:cell-diagonal}--\eqref{eq:cell-offdiagonal}; no
singular quadrature or fitted coefficient enters the tables.

\section*{AI-use disclosure}

The proof architecture and manuscript were developed with assistance from
OpenAI's GPT-5.6 Sol model (\texttt{gpt-5.6-sol}).  The model served as an
interactive research assistant for argument development and stress testing,
citation and calculation audits, preparation of reproducibility and
release-audit code, and revision.  The author directed the research, selected and
checked the final arguments, and takes responsibility for all mathematical
claims, citations, computations, and wording.

\section*{Code availability}

\begin{sloppypar}
The script \path{tools/check_asymptotics.py} reproduces
Tables~\ref{tab:mesh} and \ref{tab:clusters}.  Before solving the
eigenproblems, it checks the assembled constant-function energy identity and
prints progress after each separation.  The pinned numerical dependencies are
listed in
\path{requirements-numerics.txt}; the reference output was replayed with
Python~3.12.3, NumPy~1.26.4, and SciPy~1.11.4.  The submission source package
includes frozen copies of the scripts, pinned dependencies, and reference
output; this is the version corresponding to the manuscript.
The version tag for this manuscript is
\texttt{paper-2026-09-03-r2}.  The development repository is available at
\href{https://github.com/skypher/laplace-tunneling}%
{github.com/skypher/laplace-tunneling}.
\end{sloppypar}

\enlargethispage{3\baselineskip}
{\small
\let\standardthebibliography\thebibliography
\renewcommand{\thebibliography}[1]{%
  \standardthebibliography{#1}%
  \setlength{\itemsep}{0pt}}
\bibliographystyle{alpha}
\bibliography{references}
}

\end{document}
